\chapter{D'Artagnan's Lesson}

\lettrine{R}{aoul} did not meet with D'Artagnan the next day,  as he had hoped. He only met with Planchet, whose joy was great at seeing the  young man again, and who contrived to pay him two or three little soldierly  compliments, savouring very little of the grocer's shop. But as Raoul was  returning the next day from Vincennes at the head of fifty dragoons confided to  him by Monsieur le Prince, he perceived, in La Place Baudoyer, a man with his  nose in the air, examining a house as we examine a horse we have a fancy to  buy. This man, dressed in a citizen costume buttoned up like a military  pourpoint, a very small hat on his head, but a long shagreen-mounted sword by  his side, turned his head as soon as he heard the steps of the horses, and left  off looking at the house to look at the dragoons. It was simply M.  d'Artagnan; D'Artagnan on foot; D'Artagnan with his hands  behind him, passing a little review upon the dragoons, after having reviewed  the buildings. Not a man, not a tag, not a horse's hoof escaped his  inspection. Raoul rode at the side of his troop; D'Artagnan perceived him  the last. <Eh!> said he, <Eh! Mordioux!>

<I was not mistaken!> cried Raoul, turning his horse towards him.

<Mistaken—no! Good-day to you,> replied the ex-musketeer;  whilst Raoul eagerly pressed the hand of his old friend. <Take care,  Raoul,> said D'Artagnan, <the second horse of the fifth rank  will lose a shoe before he gets to the Pont Marie; he has only two nails left  in his off fore-foot.>

<Wait a minute, I will come back,> said Raoul.

<Can you quit your detachment?>

<The cornet is there to take my place.>

<Then you will come and dine with me?>

<Most willingly, Monsieur d'Artagnan.>

<Be quick, then; leave your horse, or make them give me one.>

<I prefer coming back on foot with you.>

Raoul hastened to give notice to the cornet, who took his post; he then  dismounted, gave his horse to one of the dragoons, and with great delight  seized the arm of M. d'Artagnan, who had watched him during all these  little evolutions with the satisfaction of a connoisseur.

<What, do you come from Vincennes?> said he.

<Yes, monsieur le chevalier.>

<And the cardinal?>

<Is very ill; it is even reported he is dead.>

<Are you on good terms with M. Fouquet?> asked D'Artagnan,  with a disdainful movement of the shoulders, proving that the death of Mazarin  did not affect him beyond measure.

<With M. Fouquet?> said Raoul; <I do not know him.>

<So much the worse! so much the worse! for a new king always seeks to get  good men in his employment.>

<Oh! the king means no harm,> replied the young man.

<I say nothing about the crown,> cried D'Artagnan; <I  am speaking of the king—the king, that is M. Fouquet, if the cardinal is  dead. You must contrive to stand well with M. Fouquet, if you do not wish to  molder away all your life as I have moldered. It is true you have, fortunately,  other protectors.>

<M. le Prince, for instance.>

<Worn out! worn out!>

<M. le Comte de la Fère?>

<Athos! Oh! that's different; yes, Athos—and if you have any  wish to make your way in England, you cannot apply to a better person; I can  even say, without too much vanity, that I myself have some credit at the court  of Charles II. There is a king—God speed him!>

<Ah!> cried Raoul, with the natural curiosity of well-born young  people, while listening to experience and courage.

<Yes, a king who amuses himself, it is true, but who has had a sword in  his hand, and can appreciate useful men. Athos is on good terms with Charles  II. Take service there, and leave these scoundrels of contractors and  farmers-general, who steal as well with French hands as others have done with  Italian hands; leave the little snivelling king, who is going to give us  another reign of Francis II. Do you know anything of history, Raoul?>

<Yes, monsieur le chevalier.>

<Do you know, then, that Francis II had always the earache?>

<No, I did not know that.>

<That Charles IV had always the headache?>

<Indeed!>

<And Henry III had always the stomach-ache?>

Raoul began to laugh.

<Well, my dear friend, Louis XIV always has the heart-ache; it is  deplorable to see a king sighing from morning till night without saying once in  the course of the day, ventre-saint-gris! corboef! or anything to rouse  one.>

<Was that the reason why you quitted the service, monsieur le  chevalier?>

<Yes.>

<But you yourself, M. d'Artagnan, are throwing the handle after the  axe; you will not make a fortune.>

<Who? I?> replied D'Artagnan, in a careless tone; <I am  settled—I had some family property.>

Raoul looked at him. The poverty of D'Artagnan was proverbial. A Gascon,  he exceeded in ill-luck all the gasconnades of France and Navarre; Raoul had a  hundred times heard Job and D'Artagnan named together, as the twins  Romulus and Remus. D'Artagnan caught Raoul's look of astonishment.

<And has not your father told you I have been in England?>

<Yes, monsieur le chevalier.>

<And that I there met with a very lucky chance?>

<No, monsieur, I did not know that.>

<Yes, a very worthy friend of mine, a great nobleman, the viceroy of  Scotland and Ireland, has endowed me with an inheritance.>

<An inheritance?>

<And a good one, too.>

<Then you are rich?>

<Bah!>

<Receive my sincere congratulation.>

<Thank you! Look, that is my house.>

<Place de Grève?>

<Yes; don't you like this quarter?>

<On the contrary, the look-out over the water is pleasant. Oh! what a  pretty old house!>

<The sign Notre Dame; it is an old cabaret, which I have transformed into  a private house in two days.>

<But the cabaret is still open?>

<Pardieu!>

<And where do you lodge, then?>

<I? I lodge with Planchet.>

<You said, just now, <This is my house.>>

<I said so, because, in fact, it is my house. I have bought it.>

<Ah!> said Raoul.

<At ten years' purchase, my dear Raoul; a superb affair; I bought  the house for thirty thousand livres; it has a garden which opens to the Rue de  la Mortillerie; the cabaret lets for a thousand livres, with the first story;  the garret, or second floor, for five hundred livres.>

<Indeed!>

<Yes, indeed.>

<Five hundred livres for a garret? Why, it is not habitable.>

<Therefore no one inhabits it; only, you see, this garret has two windows  which look out upon the Place.>

<Yes, monsieur.>

<Well, then, every time anybody is broken on the wheel or hung,  quartered, or burnt, these two windows let for twenty pistoles.>

<Oh!> said Raoul, with horror.

<It is disgusting, is it not?> said D'Artagnan.

<Oh!> repeated Raoul.

“It is disgusting, but so it is. These Parisian cockneys are sometimes  real anthropophagi. I cannot conceive how men, Christians, can make such  speculation.

<That is true.>

<As for myself,> continued D'Artagnan, <if I inhabited  that house, on days of execution I would shut it up to the very keyholes; but I  do not inhabit it.>

<And you let the garret for five hundred livres?>

<To the ferocious cabaretier, who sub-lets it. I said, then, fifteen  hundred livres.>

<The natural interest of money,> said Raoul,—<five per  cent.>

<Exactly so. I then have left the side of the house at the back,  store-rooms, and cellars, inundated every winter, two hundred livres; and the  garden, which is very fine, well planted, well shaded under the walls and the  portal of Saint-Gervais-Saint-Protais, thirteen hundred livres.>

<Thirteen hundred livres! why, that is royal!>

<This is the whole history. I strongly suspect some canon of the parish  (these canons are all rich as Croesus)—I suspect some canon of having  hired the garden to take his pleasure in. The tenant has given the name of M.  Godard. That is either a false name or a real name; if true, he is a canon; if  false, he is some unknown; but of what consequence is it to me? he always pays  in advance. I had also an idea just now, when I met you, of buying a house in  the Place Baudoyer, the back premises of which join my garden, and would make a  magnificent property. Your dragoons interrupted my calculations. But come, let  us take the Rue de la Vannerie: that will lead us straight to M.  Planchet's.> D'Artagnan mended his pace, and conducted Raoul  to Planchet's dwelling, a chamber of which the grocer had given up to his  old master. Planchet was out, but the dinner was ready. There was a remains of  military regularity and punctuality preserved in the grocer's household.  D'Artagnan returned to the subject of Raoul's future.

<Your father brings you up rather strictly?> said he.

<Justly, monsieur le chevalier.>

<Oh, yes, I know Athos is just; but close, perhaps?>

<A royal hand, Monsieur d'Artagnan.>

<Well, never want, my boy! If ever you stand in need of a few pistoles,  the old musketeer is at hand.>

<My dear Monsieur d'Artagnan!>

<Do you play a little?>

<Never.>

<Successful with the ladies, then?—Oh! my little Aramis! That, my  dear friend, costs even more than play. It is true we fight when we lose; that  is a compensation. Bah! that little sniveller, the king, makes winners give him  his revenge. What a reign! my poor Raoul, what a reign! When we think that, in  my time, the musketeers were besieged in their houses like Hector and Priam in  the city of Troy; and the women wept, and then the walls laughed, and then five  hundred beggarly fellows clapped their hands and cried, <Kill!  kill!> when not one musketeer was hurt. Mordioux! you will never see  anything like that.>

<You are very hard upon the king, my dear Monsieur d'Artagnan and  yet you scarcely know him.>

<I! Listen, Raoul. Day by day, hour by hour,—take note of my  words,—I will predict what he will do. The cardinal being dead, he will  fret; very well, that is the least silly thing he will do, particularly if he  does not shed a tear.>

<And then?>

<Why, then he will get M. Fouquet to allow him a pension, and will go and  compose verses at Fontainebleau, upon some Mancini or other, whose eyes the  queen will scratch out. She is a Spaniard, you see,—this queen of ours;  and she has, for mother-in-law, Madame Anne of Austria. I know something of the  Spaniards of the house of Austria.>

<And next?>

<Well, after having torn the silver lace from the uniforms of his Swiss,  because lace is too expensive, he will dismount his musketeers, because oats  and hay of a horse cost five sols a day.>

<Oh! do not say that.>

<Of what consequence is it to me? I am no longer a musketeer, am I? Let  them be on horseback, let them be on foot, let them carry a larding-pin, a  spit, a sword, or nothing—what is it to me?>

<My dear Monsieur d'Artagnan, I beseech you speak no more ill of  the king. I am almost in his service, and my father would be very angry with me  for having heard, even from your mouth, words injurious to his majesty.>

<Your father, eh! He is a knight in every bad cause. Pardieu! yes, your  father is a brave man, a Caesar, it is true—but a man without  perception.>

<Now, my dear chevalier,> exclaimed Raoul, laughing, <are you  going to speak ill of my father, of him you call the great Athos? Truly you are  in a bad vein to-day; riches render you as sour as poverty renders other  people.>

<Pardieu! you are right. I am a rascal and in my dotage; I am an unhappy  wretch grown old; a tent-cord untwisted, a pierced cuirass, a boot without a  sole, a spur without a rowel;—but do me the pleasure to add one  thing.>

<What is that, my dear Monsieur d'Artagnan?>

<Simply say: <Mazarin was a pitiful wretch.>>

<Perhaps he is dead.>

<More the reason—I say was; if I did not hope that he was dead, I  would entreat you to say: <Mazarin is a pitiful wretch.> Come, say  so, say so, for love of me.>

<Well, I will.>

<Say it!>

<Mazarin was a pitiful wretch,> said Raoul, smiling at the  musketeer, who roared with laughter, as in his best days.

<A moment,> said the latter; <you have spoken my first  proposition, here is the conclusion of it,—repeat, Raoul, repeat:  <But I regret Mazarin.>>

<Chevalier!>

<You will not say it? Well, then, I will say it twice for you.>

<But you would regret Mazarin?>

And they were still laughing and discussing this profession of principles, when  one of the shop-boys entered. <A letter, monsieur,> said he,  <for M. d'Artagnan.>

<Thank you; give it me,> cried the musketeer.

<The handwriting of monsieur le comte,> said Raoul.

<Yes, yes.> And D'Artagnan broke the seal.

<Dear friend,> said Athos, <a person has just been here to  beg me to seek for you, on the part of the king.>

<Seek me!> said D'Artagnan, letting the paper fall upon the  table. Raoul picked it up, and continued to read aloud:—

<Make haste. His majesty is very anxious to speak to you, and expects you  at the Louvre.>

<Expects me?> again repeated the musketeer.

<He, he, he!> laughed Raoul.

<Oh, oh!> replied D'Artagnan. <What the devil can this  mean?>  